%--------------------------------------------------------------------------------
%
%--------------------------------------------------------------------------------
\unnumberedChapter{LCS Node Firmware Design}

The previous chapters introduced the overall Layout Control System architecture,
the communication concepts and the software layers used to implement an LCS
node. The following chapters will present complete node implementations such as
the base station, block controller and other hardware modules.

Before designing individual nodes, this chapter describes the general approach
to writing firmware for an LCS node. The focus is not on the technical details of
using the individual libraries. These topics were already covered in the
previous chapters. Instead, this chapter provides guidelines for the firmware
designer and describes how the different concepts of the LCS architecture are
combined when creating a new node.

An LCS node is more than a controller connected to hardware. It represents a
logical function within the layout control system. Good firmware design
therefore starts by defining what the node represents, which capabilities it
offers and which parts of its behavior should be implemented in firmware and
which parts should be controlled through configuration.

The overall design process can be summarized as follows:

\begin{smallGapItemize}
    \item Define the purpose of the node.
    \item Identify the logical functions represented by ports.
    \item Define configurable values and attributes.
    \item Define events produced and consumed by the node.
    \item Map the logical functions to hardware resources.
    \item Implement the required callbacks and node specific behavior.
\end{smallGapItemize}

Following these steps results in firmware that is reusable, configurable and
independent from unnecessary hardware details.

%--------------------------------------------------------------------------------
\unnumberedSection{Designing the Node Concept}

The first step in developing firmware is defining the role of the node within
the layout. A node should represent a logical function rather than simply a
collection of hardware pins.

For example, a block controller is not defined by the number of input and output
pins it contains. It is defined by the fact that it manages one or more block
sections of the layout. Likewise, a turnout controller is defined by the
turnouts it controls.

The firmware designer should therefore first identify the objects that are
managed by the node. These objects are represented by ports. A port provides a
logical endpoint that can be addressed, configured and controlled independently
from the physical hardware implementation.

Typical questions when designing a new node are:

\begin{smallGapItemize}
    \item What function does this node provide in the layout?
    \item Which independent functions should become ports?
    \item Which values should be configurable?
    \item Which conditions should generate events?
    \item Which external events should this node react to?
\end{smallGapItemize}

A well-designed node exposes meaningful functions instead of exposing its
internal hardware implementation.

%--------------------------------------------------------------------------------
\unnumberedSection{Configuration versus Firmware}

A fundamental design goal of the LCS architecture is to separate fixed firmware
functionality from configurable behavior.

A firmware update should normally only be required when adding new capabilities,
correcting errors or supporting new hardware. Changes such as assigning a
different turnout address, changing event behavior or adjusting operating
parameters should be handled through configuration.

This concept is similar to the configuration variables found in DCC decoders.
The firmware provides the available capabilities, while configuration determines
how these capabilities are used in a particular installation.

The separation between firmware and configuration is one of the fundamental
design principles of the Layout Control System.

The firmware defines the capabilities of a node. Configuration defines how
these capabilities are used in a particular installation. A turnout controller,
for example, contains the capability to control a turnout, detect its state and
generate events. Which physical turnout it represents, which event activates it,
and which other nodes are informed about its state are configuration decisions.

This separation allows the same firmware image to be used in many different
applications. A firmware update is therefore only required when the capability
of the node changes, not when the layout configuration changes.

The firmware designer should therefore resist the temptation to encode layout
specific behavior directly into the application code. Whenever possible,
behavior should be represented through configurable node items, port items and
event assignments.

LCS nodes provide this flexibility through node and port items. These items
represent values that can be queried and modified by other nodes or by a
configuration system. Values that need to survive a power cycle can be stored in
non-volatile memory, while operational values can be maintained only in memory.

In addition to simple data storage, items can invoke firmware callbacks. This
allows a node to expose functions or dynamically calculated values without
requiring special communication mechanisms.

%--------------------------------------------------------------------------------
\unnumberedSection{Ports and Channels}

Ports are the logical interfaces provided by a node. While node items represent
general node attributes, ports represent independent functions that can be
addressed and controlled.

A good example is a block controller with four independent block sections. Each
block section should normally be represented by a separate port. The physical
implementation, such as input pins, current sensors or extension modules, is an
internal detail of the node.

The combination of node and port identifier forms the address of the logical
object within the LCS system. This allows other nodes to access a function
without knowing where it is physically implemented.

Channels are used internally by a port to represent individual controllable or
measurable elements. Depending on the node type, a port may provide one or more
channels. A channel may represent, for example, an individual output, sensor
value or controlled element belonging to a port.

The firmware designer should decide early which functions deserve their own
port and how the internal channels are organized.

A port should represent something that has an independent meaning in the layout,
not simply a hardware resource. A block controller with four track sections
should therefore expose four block ports. The fact that these ports are
implemented using current sensors, GPIO pins or extension boards is an internal
implementation detail.

%--------------------------------------------------------------------------------
\unnumberedSection{Event Handling}

The Layout Control System is fundamentally an event driven system. Events allow
nodes to react to changes without continuously polling other nodes.

For example, changing a turnout position can generate an event. Other nodes,
such as a control panel or display module, can subscribe to this event and
update their own state. The turnout controller does not need to know which
nodes are interested in this information.

When designing firmware, the designer should consider two directions:

\begin{smallGapItemize}
    \item Which local conditions should generate events?
    \item Which external events should this node react to?
\end{smallGapItemize}

Incoming events are mapped to ports through the event configuration table. When
a matching event is received, the callback registered for the corresponding port
is invoked.

The event mechanism is one of the most powerful configuration features of the
LCS architecture. It allows complex layout behavior to be created by
configuration instead of adding special cases to firmware.

A good example is route setting. A node containing a control button can
broadcast a ``set route'' event. All turnout ports belonging to that route can
be configured to react to this event. Each port then performs its local action
through its callback function. The individual turnouts do not need to know about
the route, and the node creating the event does not need to know where the
turnouts are located.

This event driven approach avoids unnecessary dependencies between nodes and
allows complex layout behavior to be created by configuration.

%--------------------------------------------------------------------------------
\unnumberedSection{Software Layers}

The LCS software architecture is designed so that the firmware designer can
concentrate on the unique functionality of the node while common tasks are
handled by the lower layers.

The node specific firmware communicates with the framework through defined APIs
and callback functions. In practice, a large part of firmware development
consists of implementing these callbacks and using the available driver
interfaces.

The main program of a node is intentionally simple:

\begin{smallGapEnumerate}
    \item Initialize the required libraries.
    \item Register callbacks.
    \item Perform node specific initialization.
    \item Enter the LCS processing loop.
\end{smallGapEnumerate}

Message handling, event processing, configuration management and hardware
abstraction are performed by the lower software layers.

The following chapters explain the programming model in more detail and show how
these concepts are applied in real firmware implementations.

%--------------------------------------------------------------------------------
\unnumberedSection{Periodic Tasks}

Although the LCS architecture is event driven, some hardware functions require
periodic processing. Examples are sensor monitoring, timeout handling or
checking power consumption.

The LCS core library provides periodic callbacks for this purpose. The firmware
designer implements the required processing in these callbacks rather than
creating an independent timing mechanism.

For example, a block controller may periodically read the track current
measurement. If the measured value exceeds a configured limit, the firmware can
broadcast an event and disable track power.

Direct timer programming is still possible, but should only be used when the
application requires timing characteristics beyond the capabilities provided by
the framework.

%--------------------------------------------------------------------------------
\unnumberedSection{Command Line and Display}

The LCS library already offers a command line interface. It provides commands
for managing a node and accessing its data. While this interface is not used
during normal operation, it is extremely useful during firmware development,
testing and debugging.

All commands are also available through LCS messages, allowing troubleshooting
and monitoring even when a node is already installed in the layout.

Following this philosophy, implementing node specific commands is encouraged.
Adding a new command simply requires implementing a command callback and
registering it during node initialization.

%--------------------------------------------------------------------------------
\unnumberedSection{Configuration}

Configuration of an LCS node consists of two different aspects.

The first aspect is hardware configuration. This includes mapping controller
resources, configuring extension boards and preparing the hardware interfaces.
This work is normally handled by the CDC and driver layers. The firmware designer
only needs to interact with these layers when implementing functionality beyond
the provided abstractions.

The second aspect is functional configuration. This defines how the node behaves
within the layout. The firmware designer should expose all values that may
reasonably change during installation as node or port items.

A good rule is that firmware should implement capabilities, while configuration
should define their use.

The larger node examples presented later in this book demonstrate this approach
in practice.

%--------------------------------------------------------------------------------
\unnumberedSection{Main Code}

The main code of an LCS node consists mainly of node specific initialization and
callback routines. During startup the required callbacks are registered and
control is then transferred to the LCS processing loop.

For small nodes, the initialization code and callbacks may be contained in a
single source file. Larger nodes will typically separate functionality into
multiple files or C++ classes.

One restriction of the callback mechanism should be considered when structuring
the firmware. Callbacks are function pointers and therefore cannot directly be
instance methods of a C++ object. They must either be standalone functions,
file-local functions or static class methods.

To support object-oriented firmware designs, callback registration provides a
user data parameter. This parameter can contain an object instance pointer,
allowing the callback function to forward the event to the appropriate object
instance.

%--------------------------------------------------------------------------------
\unnumberedSection{Summary}

Designing LCS node firmware starts with defining the logical purpose of the
node rather than starting with hardware details. Ports describe the capabilities
offered by the node, items provide configurable values and events allow nodes to
react to changes throughout the layout.

The different software layers introduced throughout this part of the book allow
the firmware designer to concentrate on the unique functionality of the node.
The LCS runtime manages communication, configuration and event processing, while
the CDC and driver layers isolate the hardware implementation.

In practice, node firmware consists mainly of initialization code and callback
functions that connect the node specific behavior to the LCS framework. By
carefully separating configurable behavior from fixed firmware functionality,
new hardware modules can be developed while maintaining a consistent software
architecture.

The following chapters apply these design principles to complete LCS node
implementations, where hardware design and firmware development are presented
together.